\documentclass[11pt]{article}

\usepackage[margin=0.72in]{geometry}
\usepackage{amsmath,amssymb,amsfonts}
\usepackage{array,booktabs}
\usepackage{enumitem}
\usepackage{tikz}
\usepackage{xcolor}
\usepackage{tcolorbox}
\usepackage{fancyhdr}
\usepackage{lastpage}

% =========================
% Novel Prep Brand Colors
% =========================
\definecolor{nppurple}{RGB}{98,54,255}
\definecolor{nppurpledark}{RGB}{76,42,204}
\definecolor{nplight}{RGB}{240,235,255}

% =========================
% Page Style
% =========================
\pagestyle{fancy}
\fancyhf{}
\lhead{\textcolor{nppurpledark}{\textbf{Novel Prep Learning Center}}}
\rhead{\textcolor{nppurpledark}{\textbf{SAT Module II --- Step-by-Step Solutions}}}
\cfoot{\textcolor{nppurpledark}{\thepage\ of \pageref{LastPage}}}
\setlength{\headheight}{14pt}

\setlength{\parindent}{0pt}
\setlength{\parskip}{0.55em}
\renewcommand{\arraystretch}{1.2}

% =========================
% Circle Number
% =========================
\newcommand{\circnum}[1]{%
\tikz[baseline=(char.base)]{
\node[
circle,
draw=nppurpledark,
fill=nplight,
inner sep=1.4pt,
line width=0.8pt,
text=nppurpledark
] (char) {\small\textbf{#1}};
}}

% =========================
% Boxes
% =========================
\newtcolorbox{answerbox}{
  colback=nplight,
  colframe=nppurpledark,
  boxrule=0.85pt,
  arc=2pt,
  left=8pt,right=8pt,top=8pt,bottom=8pt
}

\begin{document}

\begin{center}
    {\LARGE \textbf{SAT Module II Test}}\\[0.25em]
    {\Large \textbf{Step-by-Step Answer Key}}\\[0.25em]
    {\large Novel Prep Learning Center}
\end{center}

\begin{answerbox}
\textbf{Final Answer Sheet}

\[
\begin{array}{c|c@{\qquad}c|c@{\qquad}c|c@{\qquad}c|c}
1 & B & 2 & C & 3 & D & 4 & B\\
5 & 1680 & 6 & 700 & 7 & D & 8 & 1.1\\
9 & D & 10 & A & 11 & D & 12 & C\\
13 & D & 14 & 60000 & 15 & B & 16 & B\\
17 & -19 & 18 & 17.2 & 19 & D & 20 & D\\
21 & 12 & 22 & 9216
\end{array}
\]
\end{answerbox}

\bigskip

% =====================================================
\textbf{\circnum{1}}

A giant panda is predicted to have a mass between
\[
\frac{1}{950}
\quad\text{and}\quad
\frac{1}{850}
\]
times its mother's mass.

The mother's mass is \(95{,}950\) grams.

So the panda's mass \(p\) is between
\[
\frac{95{,}950}{950}
\quad\text{and}\quad
\frac{95{,}950}{850}.
\]

Compute:
\[
\frac{95{,}950}{950}=101
\]
and
\[
\frac{95{,}950}{850}\approx 112.88.
\]

Therefore,
\[
101\le p\le 112.89.
\]

\[
\boxed{B}
\]

\bigskip

% =====================================================
\textbf{\circnum{2}}

The function is
\[
f(x)=9000(1+0.03)^x=9000(1.03)^x.
\]

Here, \(x\) represents the number of years after the investment was made.

So \(f(2)\approx 9474.08\) means:

\[
\text{2 years after the investment was made, the CD was worth about } \$9474.08.
\]

\[
\boxed{C}
\]

\bigskip

% =====================================================
\textbf{\circnum{3}}

Hannah saves
\[
\frac15h.
\]

Wyatt saves
\[
\frac27w.
\]

Together, they save \(3270\), so
\[
\frac15h+\frac27w=3270.
\]

Multiply both sides by \(35\):
\[
35\left(\frac15h\right)+35\left(\frac27w\right)=35(3270).
\]

\[
7h+10w=114{,}450.
\]

\[
\boxed{D}
\]

\bigskip

% =====================================================
\textbf{\circnum{4}}

Data set A has \(21\) values, so the median is the \(11\)th value.

From the table:

\[
\begin{array}{c|c|c}
\text{Value} & \text{Frequency} & \text{Positions}\\
\hline
0 & 1 & 1\\
1 & 3 & 2\text{--}4\\
2 & 4 & 5\text{--}8\\
3 & 5 & 9\text{--}13
\end{array}
\]

The \(11\)th value is \(3\), so the median of data set A is
\[
3.
\]

After removing the incorrect value \(19\), data set B has \(20\) values.

The median is the average of the \(10\)th and \(11\)th values. Both are still \(3\).

So the medians are equal.

\[
\boxed{B}
\]

\bigskip

% =====================================================
\textbf{\circnum{5}}

The volume ratio is
\[
\frac{10584\pi}{392\pi}=27.
\]

For similar solids,
\[
\text{volume ratio}=(\text{scale factor})^3.
\]

So
\[
\text{scale factor}=3.
\]

Surface area scales by the square of the scale factor:
\[
3^2=9.
\]

Now find the surface area of cylinder A.

Given:
\[
V=\pi r^2h,\qquad r=7.
\]

\[
392\pi=\pi(7^2)h
\]

\[
392=49h
\]

\[
h=8.
\]

Surface area formula:
\[
SA=2\pi r^2+2\pi rh.
\]

\[
SA_A=2\pi(7^2)+2\pi(7)(8)
\]

\[
SA_A=98\pi+112\pi=210\pi.
\]

So
\[
k=210.
\]

Cylinder B has surface area
\[
9(210\pi)=1890\pi.
\]

So
\[
n=1890.
\]

Therefore,
\[
n-k=1890-210=1680.
\]

\[
\boxed{1680}
\]

\bigskip

% =====================================================
\textbf{\circnum{6}}

Let the original number of prey and predators each be \(N\).

Prey increased by \(1900\%\). This means the final number is
\[
N+19N=20N.
\]

Predators increased by \(150\%\). This means the final number is
\[
N+1.5N=2.5N.
\]

Now find how much greater the prey population is than the predator population:

\[
\frac{20N-2.5N}{2.5N}\times 100
\]

\[
=\frac{17.5N}{2.5N}\times 100
\]

\[
=7\times 100=700.
\]

\[
\boxed{700}
\]

\bigskip

% =====================================================
\textbf{\circnum{7}}

Zuri already ran \(4\) miles.

She needs to run an additional \(x\) miles.

Her goal is at least \(16\) miles.

So:
\[
4+x\ge 16.
\]

\[
\boxed{D}
\]

\bigskip

% =====================================================
\textbf{\circnum{8}}

The height function is
\[
h(t)=-16t^2+b.
\]

At \(t=0\), the object is \(19.36\) feet above the ground.

So
\[
h(0)=b=19.36.
\]

Thus
\[
h(t)=-16t^2+19.36.
\]

The object hits the ground when
\[
h(t)=0.
\]

\[
0=-16t^2+19.36
\]

\[
16t^2=19.36
\]

\[
t^2=1.21
\]

\[
t=1.1.
\]

\[
\boxed{1.1}
\]

\bigskip

% =====================================================
\textbf{\circnum{9}}

The function is
\[
f(x)=56(0.19)^x.
\]

When \(x\) increases by \(1\), the value is multiplied by
\[
0.19.
\]

That means the new value is \(19\%\) of the previous value.

So the decrease is
\[
100\%-19\%=81\%.
\]

\[
\boxed{D}
\]

\bigskip

% =====================================================
\textbf{\circnum{10}}

The equation is
\[
x^2+\sqrt{k-3}\,x+42=0.
\]

For a quadratic equation to have exactly one real solution, the discriminant must be \(0\).

\[
b^2-4ac=0.
\]

Here:
\[
a=1,\qquad b=\sqrt{k-3},\qquad c=42.
\]

Substitute:
\[
(\sqrt{k-3})^2-4(1)(42)=0.
\]

\[
k-3-168=0
\]

\[
k=171.
\]

\[
\boxed{A}
\]

\bigskip

% =====================================================
\textbf{\circnum{11}}

The circle is
\[
x^2+y^2=36.
\]

The point is \((-5,y)\), where \(y<0\).

Substitute \(x=-5\):

\[
(-5)^2+y^2=36
\]

\[
25+y^2=36
\]

\[
y^2=11.
\]

Since \(y<0\),
\[
y=-\sqrt{11}.
\]

The line is
\[
y=mx+\frac b4.
\]

Substitute \((-5,-\sqrt{11})\):

\[
-\sqrt{11}=m(-5)+\frac b4
\]

\[
-\sqrt{11}=-5m+\frac b4
\]

\[
\frac b4=5m-\sqrt{11}.
\]

Multiply by \(4\):

\[
b=20m-4\sqrt{11}.
\]

\[
\boxed{D}
\]

\bigskip

% =====================================================
\textbf{\circnum{12}}

Let \(x\) be the number of grams of the \(0.3\%\) solution.

Then the amount of \(0.15\%\) solution is
\[
80-x.
\]

Convert percentages to decimals:
\[
0.3\%=0.003,\qquad 0.15\%=0.0015.
\]

The amount of sodium chloride is:
\[
0.003x+0.0015(80-x)=0.21.
\]

Expand:
\[
0.003x+0.12-0.0015x=0.21.
\]

\[
0.0015x+0.12=0.21
\]

\[
0.0015x=0.09
\]

\[
x=60.
\]

\[
\boxed{C}
\]

\bigskip

% =====================================================
\textbf{\circnum{13}}

The equation is
\[
\frac{x}{5}+\frac{y}{9}=\frac{47}{45}.
\]

The total resistance is the sum of the individual resistances.

So each type A resistor contributes
\[
\frac15
\]
ohm.

Each type B resistor contributes
\[
\frac19
\]
ohm.

Thus:
\[
a=\frac15,\qquad b=\frac19.
\]

The positive difference is
\[
\frac15-\frac19.
\]

Use common denominator \(45\):

\[
\frac{9}{45}-\frac{5}{45}=\frac4{45}.
\]

\[
\boxed{D}
\]

\bigskip

% =====================================================
\textbf{\circnum{14}}

Object A is \(444\%\) of object B.

\[
A=4.44B.
\]

Object A is \(0.740\%\) of object C.

\[
A=0.00740C.
\]

Set the expressions equal:

\[
4.44B=0.00740C.
\]

Solve for \(C\):

\[
C=\frac{4.44}{0.00740}B.
\]

\[
C=600B.
\]

If \(C\) is \(p\%\) of \(B\), then

\[
C=\frac{p}{100}B.
\]

So

\[
\frac p{100}=600.
\]

\[
p=60000.
\]

\[
\boxed{60000}
\]

\bigskip

% =====================================================
\textbf{\circnum{15}}

“At most 71 inches” includes the first two height categories:

\[
48\text{--}59
\quad \text{and}\quad
60\text{--}71.
\]

Total number of participants at most 71 inches tall:

\[
60+180=240.
\]

Number from Group 2 at most 71 inches tall:

\[
10+75=85.
\]

Probability:

\[
\frac{85}{240}.
\]

Simplify by dividing by \(5\):

\[
\frac{85}{240}=\frac{17}{48}.
\]

\[
\boxed{B}
\]

\bigskip

% =====================================================
\textbf{\circnum{16}}

The system is:
\[
2x+3y=7
\]
and
\[
10x+15y=35.
\]

The second equation is just \(5\) times the first equation.

So the system represents the same line.

Let
\[
y=r.
\]

Substitute into
\[
2x+3y=7.
\]

\[
2x+3r=7
\]

\[
2x=7-3r
\]

\[
x=\frac{7-3r}{2}.
\]

So the point is

\[
\left(\frac{7-3r}{2},r\right)
=
\left(-\frac{3r}{2}+\frac72,r\right).
\]

\[
\boxed{B}
\]

\bigskip

% =====================================================
\textbf{\circnum{17}}

The equation is
\[
\frac{1}{cx}=\frac{x}{76}+\frac1c.
\]

Multiply both sides by \(76cx\):

\[
76=cx^2+76x.
\]

Rewrite:
\[
cx^2+76x-76=0.
\]

This quadratic has exactly one solution, so the discriminant is \(0\):

\[
b^2-4ac=0.
\]

Here:
\[
a=c,\qquad b=76,\qquad c_{\text{constant}}=-76.
\]

So:
\[
76^2-4(c)(-76)=0.
\]

\[
5776+304c=0.
\]

\[
304c=-5776.
\]

\[
c=-19.
\]

\[
\boxed{-19}
\]

\bigskip

% =====================================================
\textbf{\circnum{18}}

Convert
\[
7.7\frac{\text{meters}}{\text{second}^2}
\]
to
\[
\frac{\text{miles}}{\text{minute}^2}.
\]

Use:
\[
1\text{ mile}=1609\text{ meters}
\]
and
\[
1\text{ minute}=60\text{ seconds}.
\]

Since the time unit is squared:
\[
1\text{ minute}^2=60^2\text{ seconds}^2=3600\text{ seconds}^2.
\]

So:
\[
7.7\cdot \frac{3600}{1609}.
\]

\[
7.7\cdot \frac{3600}{1609}\approx 17.228.
\]

Rounded to the nearest tenth:

\[
\boxed{17.2}
\]

\bigskip

% =====================================================
\textbf{\circnum{19}}

One gallon covers \(170\) square feet.

The total fence area is \(w\) square feet.

To stain the fence once:

\[
\frac{w}{170}.
\]

To stain the fence twice:

\[
2\cdot \frac{w}{170}.
\]

\[
\frac{2w}{170}=\frac{w}{85}.
\]

\[
\boxed{D}
\]

\bigskip

% =====================================================
\textbf{\circnum{20}}

The function is
\[
f(x)=66b^x.
\]

When \(x\) increases by \(1\):

\[
f(x+1)=66b^{x+1}.
\]

\[
f(x+1)=b\cdot 66b^x.
\]

So the multiplier is
\[
b.
\]

If the function increases by \(c\%\), then the multiplier is

\[
1+\frac{c}{100}.
\]

Thus:

\[
b=1+\frac{c}{100}.
\]

\[
b-1=\frac{c}{100}.
\]

\[
c=100(b-1).
\]

\[
\boxed{D}
\]

\bigskip

% =====================================================
\textbf{\circnum{21}}

The equation is
\[
\frac{1}{24}x^2+\left(s-\frac{1}{24}t\right)x-st=0.
\]

For a quadratic equation
\[
ax^2+bx+c=0,
\]
the product of the solutions is
\[
\frac{c}{a}.
\]

Here:
\[
a=\frac{1}{24},
\qquad
c=-st.
\]

So the product of the solutions is

\[
\frac{-st}{\frac{1}{24}}.
\]

\[
=-24st.
\]

The problem says the product is

\[
-2kst.
\]

So:

\[
-2kst=-24st.
\]

Since \(s,t>0\), divide by \(-2st\):

\[
k=12.
\]

\[
\boxed{12}
\]

\bigskip

% =====================================================
\textbf{\circnum{22}}

A triangular pyramid has \(4\) triangular faces.

If all \(6\) edges are \(96\) cm, then each face is an equilateral triangle with side length \(96\).

Area of one equilateral triangle:

\[
A=\frac{\sqrt3}{4}s^2.
\]

Substitute \(s=96\):

\[
A=\frac{\sqrt3}{4}(96)^2.
\]

\[
96^2=9216.
\]

\[
A=\frac{9216\sqrt3}{4}=2304\sqrt3.
\]

There are \(4\) faces:

\[
SA=4(2304\sqrt3)=9216\sqrt3.
\]

The surface area is given as
\[
k\sqrt3.
\]

So

\[
k=9216.
\]

\[
\boxed{9216}
\]

\end{document}